\section{Fluxes and the complete retained size sum}\label{sec:summation}
Assume first that initial cumulants are trivial, so $p=0$.
Fix a bounded-complexity exceptional skeleton, including the first
redundant-contact slot. Let $s$ range over decorated merger pair schedules.
After spatial-envelope domination write $F_s$ for the full nonnegative
merger-flux weight, with all needed redundant-branch indicators retained.
It is zero on an unrealizable branch. Let $G_s$ be the Gaussian transverse
return demand determined solely by the formal forest before its first
redundancy. Then $\operatorname{supp}F_s\subset G_s$.
No post-return Gaussian invariance is asserted.

\begin{lemma}[Squared-flux sum]\label{lem:squared}
At fixed merger times,
\[
 \sum_s\int F_s^2\dd\Omega
       \leq C^aE^a\prod_i n_i^{a_i}.
\]
The number of ordered pair schedules is at most
$C^a\prod_i n_i^{2a_i}$.
\end{lemma}
\begin{proof}
At a valid branch node in slab $i$,
\[
 \sum_{p<q}|v_p-v_q|^2
 =n_i\sum_p|v_p|^2-\left|\sum_pv_p\right|^2\leq n_iE.
\]
The squared incoming flux is at most its squared relative speed.
Integrate the free normal and sum C/O choices, then induct backward
through the pair-choice tree. Every valid child preserves the energy
bound, including all intervening determined redundant scatterings.
Invalid children remain zero; adding other pair speeds only enlarges
the estimate. Thus no fictitious velocity is needed on an invalid branch.
The schedule count follows by permitting every unordered pair at every
merger, ignoring forest restrictions only in this numerical upper bound.
\end{proof}

Apply Cauchy--Schwarz jointly in schedules and the product parameter
measure $\nu$:
\begin{equation}\label{eq:jointCS}
 \sum_s\int F_s\dd\nu
 \leq\left(\sum_s\int F_s^2\dd\nu\right)^{1/2}
       \left(\sum_s\int\ind_{G_s}\dd\nu\right)^{1/2}.
\end{equation}
For $m$ birth Gaussians and $a=m-1\leq n$,
\[
 \E_\gamma E^a=(2\sigma^2)^a
       \frac{\Gamma(3m/2+a)}{\Gamma(3m/2)}
       \leq(C_bn)^a .
\]
In the nonneighbor case include \emph{all} free merger times in $\nu$.
With $V_t=\prod_i\tau^{a_i}/a_i!$ and
$B_\eps=1+\logp(\sigma\tau/\eps)$, the two factors in \eqref{eq:jointCS},
before taking square roots, are bounded by
\[
 C^{m+a}V_t n^a\prod_i n_i^{a_i},
 \qquad
 C^aV_t\prod_i n_i^{2a_i}\frac{Cn\eps B_\eps}{\sigma\tau}.
\]
Their product restores the full time volume and gives
\begin{equation}\label{eq:preStirling}
 C^{n+a}V_t n^{a/2}\prod_i n_i^{3a_i/2}
              \sqrt{\frac{n\eps B_\eps}{\sigma\tau}} .
\end{equation}
Only the second factor uses the Gaussian forest geometry. In particular
later merger weights and nonlinear determined returns stay in the first.

In the repeated-pair sector, refine the exceptional skeleton by the
earlier pair-contact slot and the first later C endpoint-update slot.
These at most $a^2\leq n^2$ choices are absorbed in the fixed-complexity
polynomial. Within each refined class $t_0$, $t_1$, and the elementary
gap selected below are common to all schedules. Classes with no such
C update vanish by \Cref{lem:repeated}.
We may now apply \eqref{eq:jointCS} at fixed times
and use \Cref{lem:repeated}. The resulting factor is
$\min(1,\sqrt{\eps/(\sigma\ell)})$.
If the earlier edge lies in slab $i$, choose either its next same-slab
merger gap or the terminal gap to that slab's endpoint. In both cases
the chosen gap $\Delta$ is at most $\ell$. Under the normalized ordered
slab simplex,
\[
 \E\Delta^{-1/2}
 =\tau^{-1/2}
   \frac{\Gamma(1/2)\Gamma(a_i+1)}{\Gamma(a_i+1/2)}
 \leq C\sqrt{(1+a_i)/\tau}.
\]
Thus \eqref{eq:preStirling} holds without $B_\eps$ in that sector.

\begin{proof}[Proof of \Cref{thm:cycles}]
Insert the fresh-label denominator $\prod_i n_i!$ and the fixed
complexity polynomial. The source's elementary Stirling calculation,
with $\delta_i=n_i-a_i\geq0$, gives
\[
 \frac{n_i^{3a_i/2}}{n_i!a_i!}
 \leq e^{2n_i}(1+n_i)^{\delta_i}n_i^{-n_i/2},
 \qquad
 \prod_i n_i^{n_i}\geq(n/L)^n .
\]
Because $\sum_i\delta_i=O_K(1)$, \eqref{eq:preStirling} is bounded by
\[
 C_{K,h,T}\mu^{1-h}r_\eps
       (1+n)^{d_K}(C_b\tau\sqrt L)^n .
\]
The activity factor is exactly $\mu^{m-h}\eps^{2a}=\mu^{1-h}$.
The number of slab-size vectors of total $n$ is polynomial for fixed $L$.
Choose $L$ once so that $C_bT/\sqrt L<1$. The size series converges.
The repeated sector gives $\sqrt\eps$ in place of $r_\eps$.
The first redundant pair is either an earlier edge or a nonneighbor,
so these cases exhaust the $p=0$ cyclic class.
\end{proof}

Physical initial-link sectors separately have $\eps^p$, $p\geq1$,
and the original summable majorant bounds them by
$C_K\mu^{1-h}\eps$. Nonphysical initial tails and high complexity
are treated by their actual arbitrary-power estimates in
\Cref{prop:imports}; they are not made independent by the new geometry.

\begin{remark}
Uniform $C^n$ flux moments for each fixed particle tree would be false.
For collinear C contacts $(1,2),(2,3),\ldots,(n-1,n)$, the same ancestral
normal scalar is transported along the chain. Its flux product is
$\prod_{j=2}^n(X_1-X_j)_+$ with independent Gaussian $X_j$.
The event $X_1/\sigma\in[2\sqrt{n-1},2\sqrt{n-1}+1]$ and
$|X_j|/\sigma\leq1$ gives expectation at least
$c^n\sigma^{n-1}(n-1)^{(n-1)/2}$.
The particle tree has degree two. Our all-schedule squared-flux proof
preserves, and then cancels, this Gaussian growth through the two
factorial denominators.
\end{remark}
